% TOP_PROBLEM: {"id":30007756,"problem_number":"LOCAL-30007756","title":"Corporate-tax incidence","category_id":21,"category":{"id":21,"name":"economics","display_name":"Economics","description":"Open economic theory statements, published research agendas, and proposed scoped specifications; question provenance and historical priority are qualified per record.","slug":"economics","order_index":21,"created_at":"2026-10-08T00:00:00Z"},"status":"research_question","external_url":null,"legacy_ids":[],"published":false,"statement":"In the explicitly specified setting: How does the corporate-tax burden divide among workers, shareholders, consumers, and landowners across markets and time horizons? The mathematical statement above fixes the target, assumptions and scope of this catalogue formulation.","economics":{"original_id":245,"field":"Taxation and social insurance","kind":"identification","formalization_basis":"adapted_research_question","source_status":"research_agenda","priority_status":"earliest_found","priority_note":"Earliest explicit agenda verified in this audit, not a claim of first historical publication. The mathematical scope is authored for this catalogue.","earliest_verified_reference":{"authors":["Alan J. Auerbach"],"title":"Who Bears the Corporate Tax? A Review of What We Know","year":2006,"url":"https://www.journals.uchicago.edu/doi/10.1086/tpe.20.20061903","doi":"10.1086/tpe.20.20061903","locator":"Tax Policy and the Economy 20, 1–40; publisher abstract","evidence_summary":"Explicitly asks who bears the corporate income tax and describes the answer as elusive, including timing and shareholder burdens. Earlier 2005 working-paper and 1962 Harberger antecedents exist."},"current_reference":{"authors":["Alan J. Auerbach"],"title":"Who Bears the Corporate Tax? A Review of What We Know","year":2006,"url":"https://www.journals.uchicago.edu/doi/10.1086/tpe.20.20061903","doi":"10.1086/tpe.20.20061903","locator":"Tax Policy and the Economy 20, 1–40; publisher abstract","evidence_summary":"Explicitly asks who bears the corporate income tax and describes the answer as elusive, including timing and shareholder burdens. Earlier 2005 working-paper and 1962 Harberger antecedents exist."},"formalization_note":"Catalogue-authored scoped formalization. Population, instruments, welfare objective and identification restrictions must be instantiated before claiming a resolution; source authors are not credited with these added assumptions."},"statusQualification":"Published question or research agenda verified at the cited locator. The mathematical specification is adapted for this catalogue and includes additional assumptions; source publication does not establish first-ever priority or certify this exact model remains unsolved. Catalogue-authored scoped formalization. Population, instruments, welfare objective and identification restrictions must be instantiated before claiming a resolution; source authors are not credited with these added assumptions.","statusReviewedAt":"2026-10-08"}
% FIRST_POSTED: 2026-10-08
% ENTRY_KIND: statement_only
% !TeX program = lualatex
\documentclass[11pt]{article}
\usepackage{fontspec}
\usepackage[a4paper,margin=25mm]{geometry}
\usepackage{amsmath,amssymb,mathtools}
\usepackage{xurl}
\usepackage[hidelinks]{hyperref}
\newcommand{\sref}[2]{\hyperlink{source:#1:#2}{[S#2]}}
\setlength{\emergencystretch}{3em}
\begin{document}
% problemId:30007756
\section*{Corporate-tax incidence}\label{30007756}
\subsection{Definitions and mathematical statement}
Specify a corporate-tax reform \(p_0\to p_1\), economy, affected firms and horizon \(H\), holding the destination of incremental public revenue fixed. For households \(i\), let \(U_i^H(p,x)\) be lifetime utility under policy \(p\) with a date-zero transfer \(x\), in a fixed numeraire. Define compensating variation by \(CV_i^H=\inf\{x\in\mathbb R:U_i^H(p_1,x)\geq U_i^H(p_0,0)\}\), allowing wages, consumer prices, rents and asset values to adjust. Assume this finite amount is well defined by monotonicity and continuity in \(x\). Define government incremental net receipts \(\Delta R_H\neq0\). For prespecified mutually exclusive household groups \(g\), the incidence measure is
\[\iota_g(H)=\frac{\sum_{i\in g}CV_i^H}{\Delta R_H}.\]
For channels rather than people, specify an attribution rule using a structural equilibrium model: alter labor income, dividends and capital gains, consumer prices and land rents individually and average incremental effects over every channel ordering. This fixes a Shapley decomposition and avoids double-counting workers who also own shares. Estimate or sharply bound group incidence and channel contributions at short and long horizons. The model must specify capital mobility, market power, investment adjustment and expectations before reform; the empirical design must separate policy shocks from correlated changes. Report both revenue-transfer incidence and excess burden, so channel shares need not sum to one without an explicit accounting identity. Determine which cross-market data identify the allocation of losses and how results change across tax bases.

\par\textbf{Formulation and scope.} Catalogue-authored scoped formalization. Population, instruments, welfare objective and identification restrictions must be instantiated before claiming a resolution; source authors are not credited with these added assumptions.

\par\textbf{Open-question provenance.} An explicit question or research agenda was verified at the source locator. This does not by itself certify that every later version remains unresolved \sref{30007756}{1}.

\par\textbf{Historical priority.} Earliest explicit agenda verified in this audit, not a claim of first historical publication. The mathematical scope is authored for this catalogue.

\subsection{Short English statement}
In the explicitly specified setting: How does the corporate-tax burden divide among workers, shareholders, consumers, and landowners across markets and time horizons? The mathematical statement above fixes the target, assumptions and scope of this catalogue formulation.

\subsection{Sources}
\begin{itemize}\raggedright
\item[\textnormal{[S1]}]\hypertarget{source:30007756:1}{} Alan J. Auerbach. 2006. Who Bears the Corporate Tax? A Review of What We Know. Tax Policy and the Economy 20, 1–40; publisher abstract.
\url{https://www.journals.uchicago.edu/doi/10.1086/tpe.20.20061903}.
DOI: 10.1086/tpe.20.20061903.
{\small\textit{Earliest verified question/agenda reference; historical priority qualified below. Explicitly asks who bears the corporate income tax and describes the answer as elusive, including timing and shareholder burdens. Earlier 2005 working-paper and 1962 Harberger antecedents exist.}}
\end{itemize}
\end{document}
